% Transcription — fonds Alexandre Grothendieck, shelfmark 154, batch 3 (pages 41-60).
% Established from the facsimile archives/batches/154/batch-03.pdf, cut from a
% copy of 154.pdf fetched through the project's own relay
% (https://grothendieck-relay.onrender.com/source/154.pdf); PDF page k =
% archivists' page 40+k, no cover sheet. Per the skill transcribe-grothendieck.
% The folder is 175 pages in nine batches (1-20, ..., 161-175), transcribed in
% parallel; batch 2 was read first, since p. 41 continues its last sentence.
% Notation follows batches 1-2 and the transcription of folder 80.
%
% Pass: Opus 5.5 (claude-opus-5-5), 2026-09-26 — first pass, unchecked against the pages by a human.
%
% Dating and title. \dating is the folder's own date, copied verbatim from
% src/content/catalogue.ts; the folder belongs to the inventory group
% « Espaces stratifiés » (150-151), but since it has a date of its own the
% group's range is not added. Title from the catalogue, trailing full stop
% dropped; « [Système de pseudo-droites] » is bracketed there because the
% archivists supplied it.
%
% Pages skipped: none. Pages summarised: none. Pages given only as figures
% described in French \note{}s: 52 (drawings, no text of his); 56 is mostly
% drawings, with one formula transcribed. Page 55 is a page of pencil
% calculations, set formula by formula. All other pages (41-51, 53-54,
% 57-60) are transcribed from his hand, fast drafting with many \ill{}.
%
% His pagination and dates. No page number of his is visible in this batch.
% Dates in his hand: « 31.12.83 » at the head of p. 43 (which opens a new
% redaction headed « Déploiement d'un système de pseudo-droites »);
% « 7.1.1984 » on p. 47; « 2.1.84 » on p. 48 (first digit read by
% comparison with p. 50) and « 2.1.84. » on p. 50; « 3.1.84 » on p. 57 (last
% digit probable). The leaves are therefore not in date order (p. 47 is
% dated after pp. 48-57). Lettered runs: a)-n) on pp. 43-46 (with f'), the
% cases a) b) c) inside f), (h) 1)-4) on p. 44, (i) 1)-4), (j), k) on p. 45,
% l) m) n) on p. 46); a) b) (c) on pp. 47, 49; a) b) c) on p. 50; d) with
% a) b) and « En résumé » a) b) c) on p. 51; the Proposition a) b) c) and
% cases 1°-3° on p. 53; « 4.1. » on p. 60. The printed job banners of the
% computer-listing paper (1982 machine dates) are paper, not his dates, and
% are not recorded.
%
% What the batch is about.
%   41-42  (end of the 25.12.83 run of batch 2) The case L = D_i in Sigma:
%          Delta_L with n-1 chords, the specialisation L -> D, gluing Delta_L
%          to Delta_D by the canonical isomorphism, again the opposite of the
%          tautological isomorphism on orientations; generisation arcs of D
%          overlap, whence the idea of « multiplying » D per vertex s, so that
%          the exceptional generisation arcs bound bigons or polygons without
%          inertia arcs (« un triomphe »): the compact surface is completed by
%          bigons and squares.
%   43-46  31.12.83, « Déploiement »: axioms for the unfolding surface
%          X(Sigma) with two graphs K, L: vertices of K of order 4, ordinary
%          and exceptional edges, K_0 a union of circles indexed by I, critical
%          points, multiplicities nu(s), the count 2n of points of L on each
%          face boundary; (h) indexation of faces, edges, ordinary and
%          exceptional vertices of K by relative positions, immediate
%          specialisations and flags; (i) the structural map phi : X -> X
%          (foldings, double foldings, L = phi^{-1}(union D_i)); (j) the face
%          F' and the affine situation; k) stable / unstable exceptional edges;
%          l)-n) symmetry law around an exceptional edge, index of exceptional
%          edges, reconstruction of Sigma from K and the numerical data;
%          involution sigma_a on a tubular neighbourhood.
%   47-54  Structure of K: red and black (orange) vertices S'_K, S''_K,
%          signs epsilon_F(a) and their rules, the symmetry bijections
%          sigma_a^F : Pol(F) -> Pol(F') (pp. 47, 49-51), two drafts of the
%          placing of new vertices (pp. 48, 50, both 2.1.84), critical points
%          C_K with weights nu(s), the relation between critical points on F
%          and on sigma_a(F); Proposition (p. 53): the cases (3,3), (3,4),
%          (3,5); Corollaries (p. 54): the signs for one face determine them
%          for all; faces of supercritical positions still to be examined.
%   55-56  Pencil calculations: l_s = n - nu_s, l_a, chi_{2n} = n - 1; the
%          parallel transport lambda(t) = t - (n-1) along a half-turn,
%          composites of reflections sigma_i(t) = l_i - t.
%   57-60  3.1.84: parallel transport along a half-turn round a
%          supersingular position, lambda(t) = t - (n-1); image of
%          pi_1(X) -> D_2n (gcd of n-1 and 2n), of pi_1(X~) -> D_2n; parity of
%          the number of elementary operations (« merveilleux, c'est toujours
%          pair ! »), corrected to 2(...) + nu(s); surjectivity of
%          pi_1(X) -> D_2n; 4.1: transport Pol(c) along combinatorial paths,
%          the characters chi_*, chi_s, chi_a of pi_1(X, F_D).
%
% Notation fixed once for the batch, following batches 1-2: \underline{\Sigma}
% for his doubly underlined Sigma; \mathcal{X} for his script X (the unfolding
% surface), \widetilde{\mathcal{X}} for its covering; \Delta_L, \widetilde{L};
% \operatorname{Pol}(D), \operatorname{Pol}(F); \underline{\omega} for his
% underlined omega; \mathbb{D}_{2n} for the dihedral group (his double-struck
% D); S_K, S'_K, S''_K, C_K; \varepsilon_F(a); \sigma_a^F.
%
% Readings to check first: pp. 43-45 (the axioms, dense with interlinear
% additions); p. 51 (d) and « En résumé »); the couples of p. 53; the tilde
% marks distinguishing the two coverings on p. 58; the ends of pp. 57-59.
\documentclass[11pt,a4paper]{article}
% Le préambule commun à toutes les transcriptions du fonds Grothendieck.
%
% Il tient en une idée : **l'incertitude se compose, elle ne se résout pas.**
% Une transcription qui lisse un mot douteux a détruit la seule chose pour
% laquelle on la faisait — un lecteur doit pouvoir savoir, à chaque endroit,
% ce qui était sur le feuillet et ce qui a été ajouté. Les six macros qui
% suivent sont donc l'appareil critique, et non de la décoration.
%
% Les mêmes commandes sont reconnues par `scripts/render.mjs`, qui produit la
% vue de lecture du site. Une macro ajoutée ici doit l'être aussi là-bas,
% faute de quoi le rendu échouera bruyamment — ce qui est voulu : un rendu
% silencieusement faux est pire qu'un rendu absent.

\ProvidesPackage{grothendieck}[2026/08/08 Transcriptions du fonds Grothendieck]

\usepackage{amsmath,amssymb,amsthm}
\usepackage{mathrsfs}
\usepackage{stmaryrd}
% Les diagrammes commutatifs sont partout dans ce fonds : c'est la langue
% même de la géométrie algébrique de Grothendieck, et une transcription qui
% les rendrait en prose ne transcrirait rien.
\usepackage{tikz-cd}
% Français par défaut — c'est la langue des feuillets — mais l'anglais est
% chargé aussi : la traduction et le résumé sont des documents anglais, et la
% typographie française (espaces avant les deux-points) y serait une faute.
% Ces documents basculent par \selectlanguage{english} en tête de corps.
\usepackage[english,french]{babel}
\usepackage{fontspec}
\usepackage{geometry}
\geometry{a4paper,margin=25mm,marginparwidth=28mm}
\usepackage{marginnote}
\usepackage{xcolor}
% ulem, not soul. `soul`'s \st and \ul are built on a character-by-character
% scan that XeLaTeX's font machinery breaks: the document fails with an
% undefined control sequence rather than degrading. ulem does the same two jobs
% and is Unicode-engine safe. `normalem` stops it hijacking \emph, which the
% transcriptions use for Grothendieck's own emphasis.
\usepackage[normalem]{ulem}
\usepackage{hyperref}
\hypersetup{colorlinks=true,linkcolor=black,urlcolor=black,pdfborder={0 0 0}}

% --- Métadonnées du lot -----------------------------------------------------
% Reprises telles quelles par le script de rendu, qui les affiche en tête de
% la vue de lecture. Elles fixent ce que le fichier prétend couvrir : sans
% elles, une transcription flotte sans point d'attache dans un fonds de seize
% mille feuillets.

\newcommand{\folder}[1]{\gdef\grothFolder{#1}}
\newcommand{\batch}[1]{\gdef\grothBatch{#1}}
\newcommand{\leaves}[2]{\gdef\grothFirst{#1}\gdef\grothLast{#2}}
\newcommand{\foldertitle}[1]{\gdef\grothFoldertitle{#1}}
\gdef\grothFolder{?}\gdef\grothBatch{?}\gdef\grothFirst{?}\gdef\grothLast{?}\gdef\grothFoldertitle{}

% --- Le résumé --------------------------------------------------------------
%
% Il ouvre la lecture modernisée, sous le titre « Résumé », et il est composé
% en petit. Ce n'est pas un ornement : c'est le seul passage du document qui
% ne soit pas des mathématiques, et un lecteur doit pouvoir le reconnaître
% avant de l'avoir lu — puis le sauter s'il connaît déjà le sujet, ou s'y
% arrêter s'il ne le connaît pas. Le filet qui le suit marque la reprise du
% corps.

\newenvironment{resume}
  {\small\setlength{\parskip}{0.45em}}
  {\par\medskip\hrule\medskip}

% --- Mots-clés ---------------------------------------------------------------
%
% En anglais, en clôture du résumé de la lecture modernisée : le vocabulaire
% moderne sous lequel chercher ce que les feuillets construisent. Le manifeste
% du site (`npm run manifest`) les extrait pour étiqueter la cote entière —
% cette ligne est la source unique des tags, il n'y a pas de fichier à côté.
\newcommand{\keywords}[1]{\par\smallskip\noindent
  {\footnotesize\textsc{Keywords} --- \textit{#1}}\par}

% --- L'appareil critique ----------------------------------------------------

% Le numéro de feuillet, en marge. C'est l'ancre qui permet de régler un
% désaccord sur une formule en regardant *un* feuillet plutôt qu'une chemise
% de six cents. Le site s'en sert aussi pour tourner les pages du fac-similé
% quand on fait défiler la transcription.
\newcommand{\leaf}[1]{%
  \marginnote{\footnotesize\color{gray!70}\textsf{#1}}%
  \label{leaf:#1}\ignorespaces}

% Illisible : on ne devine pas. Un mot inventé qui se lit comme les autres est
% le pire résultat possible.
\newcommand{\ill}{\textcolor{red!60!black}{[\,\ldots\,]}}

% Lecture douteuse : le mot est donné, et signalé comme incertain.
\newcommand{\uncertain}[1]{\uline{#1}}

% Ajout de l'éditeur — une lettre manquante, un mot sous-entendu.
\newcommand{\add}[1]{\textcolor{blue!50!black}{[#1]}}

% Biffé par Grothendieck lui-même. Ce qu'il a rayé fait partie du document :
% une rature dit souvent où la pensée a changé de direction.
\newcommand{\struck}[1]{\sout{#1}}

% Note du transcripteur, sur l'état matériel du feuillet.
\newcommand{\note}[1]{\par\smallskip{\small\color{gray!60!black}\textit{[#1]}}\par\smallskip}

% Note marginale de Grothendieck, distincte du corps du texte.
\newcommand{\marginal}[1]{\par\smallskip\noindent\hspace{1em}%
  {\small\color{gray!60!black}#1}\par\smallskip}

% --- Titre ------------------------------------------------------------------

\AtBeginDocument{%
  \begin{center}
    {\large\bfseries Fonds Alexandre Grothendieck --- cote n\textsuperscript{o}~\grothFolder}\\[2pt]
    {\itshape\grothFoldertitle}\\[4pt]
    {\small feuillets \grothFirst--\grothLast\ (lot \grothBatch)}\\[2pt]
    {\footnotesize\color{gray!60!black}Transcription établie d'après le fac-similé numérisé
     par l'Université de Montpellier.\\
     Les lectures incertaines sont soulignées, les illisibles notées [\,\ldots\,],
     les ajouts entre crochets.}
  \end{center}
  \vspace{1em}}

\endinput


\folder{154}
\batch{3}
\pages{41}{60}
\dating{1983-1984}
\watermark{Édition de démonstration}
\foldertitle{[Système de pseudo-droites] : notes manuscrites (1983-1984, s.d.)}

\begin{document}

\page{41}
on trouve un syst.\ de $n-1$ (au lieu de \uncertain{$n+1$}) cordes dans ce disque, correspondant aux éléments de $\Sigma \smallsetminus \lbrace D_i \rbrace$. Il lui correspond un polygone combinatoire \add{d'ordre $2(n-1)$}, qui \uncertain{s'appuie} sur le polygone topologique \struck{\ill{}} \struck{\ill{}} $\Delta_L$.
\note{la page continue la fin de la page 40 (lot 2) : « pour une $L = D_i \in \Sigma$, on définit encore $\Delta_L$, en considérant $D_i$ comme p.r.\ de $\underline{\Sigma} = (X, \Sigma \smallsetminus \lbrace D_i \rbrace)$ ». « d'ordre $2(n-1)$ » est ajouté au-dessus de la ligne, avec un renvoi ; le « $n+1$ » de la parenthèse est d'une lecture \uncertain{douteuse}}

L'opération de spécialisation correspondante \add{sur $L$, \ldots} : un arc de spécialisation \ill{} à un couloir de $\Sigma \cup \lbrace L \rbrace$, de murs $(L, D_{a})$, est $L \to D$ --- sur $\Delta_L$, elle consiste à \struck{\ill{}} « tirer » l'arc de spécialisation au-dessus du couloir entre cet arc et $D_{a}$, par dessus $D$, de façon à éliminer \ill{} \ill{} $D$ (dans un terme \ill{} un polygone top.\ et un syst.\ de cordes de card $(n-1)$ inscrit dedans), mais la configuration obtenue en fait n'est autre que celle de $D$, $\Delta_D$, à isom.\ canonique près. Il est entendu que l'arc \add{\uncertain{à côté}} \uncertain{tiré} jusqu'en $y$, $y'$, devenant adjacents à $x$, $x'$, le nouvel arc qui est ainsi formé, en pointillés bleus, est \struck{\ill{}} canoniquement isomorphe (en tant que pseudo-segment topologique, ou combinatoire) à $a$. Il est donc justifié de procéder comme on l'a fait, et de recoller $\Delta_L$ à $\Delta_D$ par cet iso.\ canonique, qui donne lieu à un isom.\ de transition
\[
\begin{array}{ccc}
\omega(\Delta_L) & \longrightarrow & \omega(\Delta_D) \\
\wr & & \wr \\
\omega(L) & & \omega(D)
\end{array}
\]
qui est à nouveau l'\emph{opposé} de l'isom.\ tautologique.
\note{l'indice de « $D_{a}$ » (deux fois) est \uncertain{peu net}. Figure au bas de la colonne de gauche : un disque au bord vert (clair et foncé), traversé de cordes orange ; une corde verticale en pointillé vert et bleu joint en haut les points marqués $y$, $x$ à ceux marqués $y'$, $x'$ en bas ; la moitié gauche du disque est hachurée au crayon, la moitié droite striée de traits noirs horizontaux ; $D$ est marqué deux fois près de la corde, $a$ à droite sur le bord}

L'arc correspondant sur $\Delta_D$ est « un » arc de générisation. Chaque sommet de $D$ détermine deux arcs de générisation. L'ennui ici, c'est que \uncertain{ces arcs} de générisation empiètent les uns sur les autres !

\page{42}
Pour pallier à cet inconvénient, on pourrait songer à « multiplier » chaque $D \in \Sigma$, en autant d'exemplaires qu'il y a de sommets, de sorte que \struck{chaque} pour un $(D, s)$ donné (« éliminant le contact » de $\underline{\Sigma}$), il y aurait exactement deux \emph{arcs de} \struck{spécialisation} \add{\emph{générisation}} \add{(\emph{exceptionnelles})} correspondants sur $\Delta_D$, qui n'empiètent pas. Si le sommet $s$ est simple, alors les \struck{deux} arcs de générisation exceptionnelles sont les deux côtés d'une structure de bigône topologique, \add{\uncertain{deux}} sinon, ces arcs \add{(ordinaires)}, plus les \add{deux} arcs de générisation \add{correspondants à $s$}, forment une structure de polygone topologique \uncertain{couvrant} sur $\Delta_D$. \struck{Dans} Dans les deux cas, il n'y a pas d'arc d'inertie --- \uncertain{ce} qui est un triomphe ! Bon, on a \add{fait} ce qu'il faut pour compléter une surface compacte, en recollant à la précédente des bigônes ou carrés. \add{Et le type de graphe,} \struck{Mais il faudrait déterminer} \struck{l'iso} remplacer le graphe (orienté) précédent, par celui déduit en remplaçant chaque sommet « exceptionnel » par l'ensemble des \add{\uncertain{diagonales}} \struck{sommets} du polygone local
\note{« arcs de » est souligné ; « générisation » est écrit au-dessus de « spécialisation » biffé, souligné, et surmonté de « (exceptionnelles) » également souligné. Le « deux » ajouté après « bigône topologique, » est \uncertain{peu net}. « Et le type de graphe, » est écrit au-dessus de la ligne biffée « Mais il faudrait déterminer », sans que la syntaxe se raccorde ; la page s'arrête sur « polygone local », le bas de la colonne de droite étant blanc}

\page{43}
\section{Déploiement d'un système de pseudo-droites}
\note{« 31.12.83 » est écrit en tête de la page, à gauche du titre : date de sa main (31 décembre 1983). Le titre « Déploiement » est souligné ; il ouvre une nouvelle rédaction, qui ne porte pas de numéro de page}

\emph{Déploiement} d'un syst.\ de pseudo-droites $\underline{\Sigma} = (X, \Sigma)$, $\Sigma = (D_i)_{i \in I}$ : surface compacte $\mathcal{X} = \mathcal{X}(\underline{\Sigma})$, munie de deux sous-ensembles $K$, $L$ tels que $K$, $L$ découpent cellulairement $\mathcal{X}$.

a) Tous les pts de $K$ sont d'\emph{ordre $4$}. (\ill{} déploiement $K' \to K$ qui est \ill{} \uncertain{topologique})

\struck{a} Tous les pts de $L$ sont d'ordre pair, les \ill{} arêtes qui \ill{} \add{sont distinctes (\ill{} \ill{} \uncertain{boucles} \ill{})}.

b) \struck{Pour} toute arête de $K$ contient au plus un sommet de $L$.

c) La réunion \add{$K_0$} des adh.\ des arêtes de $K$ qui ne contiennent pas de sommet \add{(de $L$ (\emph{arêtes exceptionnelles} de $K$))} est une réunion finie de cercles, \struck{indexés par $I$}
\note{« arêtes exceptionnelles de $K$ » est souligné dans l'ajout interlinéaire}

\struck{d) Soit $K_1$ la réunion des arêtes \add{non} exceptionnelles} [Il revient au même de dire que \uncertain{les} pts de $K$ \uncertain{« voisins »} à une arête exceptionnelle sont exceptionnels, et que \struck{les points} \add{pour tout} sommet de $K$, il existe au moins une arête issue de ce sommet qui \struck{\ill{}} est « ordinaire ».
\note{le crochet ouvert n'est pas refermé ; la lecture de cette phrase interlinéaire, écrite au-dessus du « d) » biffé, est \uncertain{incertaine} dans son début}

d) \struck{Notation} Soit $K_1$ la réunion des adh.\ des arêtes \add{ordinaires} \struck{\ill{}} de sorte que \add{$K = K_0 \cup K_1$} : $K_1$ est une \struck{réunion de} \ill{} courbes et formons \ill{} $K'$ (i.e.\ réunion de comp.\ \ill{}). Alors \add{\ill{}} on \uncertain{se donne} une section au-dessus des fibres des bandes pour l'immersion $K'_1 \hookrightarrow \mathcal{X}$.
\note{« d) » est entouré. L'ajout au-dessus de « $K_1$ » (« $K = \ldots$ ») est d'une lecture \uncertain{douteuse}}
\marginal{cette donnée résulte en fait des autres, \uncertain{à savoir} de la \uncertain{position} \ill{} ; cf.\ f) plus bas}
\note{note écrite en biais dans la marge de gauche, à la hauteur de d) et e)}

e) Soit $s$ un sommet exceptionnel de $K$ (i.e.\ \struck{\ill{}} sommet incident \struck{\ill{}} à une arête exceptionnelle). Alors sur \struck{\ill{}} une arête non exceptionnelle \add{(pt critique de $K$)} \struck{issue} incidente à $s$, entre $s$ et le pr\supplied{emier} sommet de $L$, \ill{} \ill{} \ill{} \add{multiplicité des arêtes \add{ord.}\ issues de $s$}. De plus, les \ill{} pts de $L$ \add{\ill{} contigus à $K$ sur les deux} \ill{} \ill{} \ill{} \ill{}
\note{figure dans le texte de e) : un segment vert foncé sur $K$, ses deux extrémités relevées en vert clair, l'une marquée d'un point noir et d'une étoile de traits orange ; trois traits orange en pointillé traversent le segment}

f) Pour tt $x \in L \cap K$, \ill{} \ill{} branche de $L$ en $x$ \ill{} \ill{} branche de $K$ en $x$. Si $x$ est un sommet de $K$, \ill{} $x$ est un pt d'ordre $4$ \add{\uncertain{évident} ?} de $L$, et \struck{\ill{}} \ill{} \ill{} \ill{} \ill{} local de $K$ en $x$ \struck{\ill{}} \add{dans} \ill{} \ill{} \ill{} une arête de $L$ \ill{} en $x$. [Donc trois cas de figure pour un $x \in K \cap L$ :

a) $x$ \add{$\notin S(L)$} \struck{non sommet de $L$}, donc pt d'indice $2$ de $L$, alors $x \notin S(K)$ ;

b) $x \in S(L) \cap \complement S(K)$ ;

c) $x \in S(L) \cap S(K)$.
\note{figures dans la colonne de droite, marquées a), b), c) : a) une droite verte traversée par un trait orange ; b) un segment vert aux extrémités fourchues, traversé en son milieu par une étoile de traits orange ; c) une croix verte traversée d'une croix orange}

f$'$) Un sommet exceptionnel $s$ de $K$ n'est pas sommet de $L$, i.e.\ $s \notin L$.

g) \struck{Un sommet} Pour un sommet \struck{\ill{}} \add{ordinaire} $s$ de $K$, sur les $4$ arêtes \add{(ordinaires)} de $K$ qui en sortent, le nb d'él.\ \struck{\ill{} \ill{}} \add{de} $L$ \struck{(\ill{} des} entre $s$ et l'unique sommet de $L$ sur l'arête \struck{\ill{} \ill{} \ill{}} \emph{est} indépendant de l'arête, soit $\nu(s)$.

\[
\varepsilon(s) = \left\lbrace \begin{array}{ll} +1 & \text{si } s \in L \\ 0 & \text{si } s \notin L \end{array} \right.
\]
\note{cette définition, introduite par « \struck{Soit} », est biffée de plusieurs traits sur la page}

Pour toute arête ordinaire, soit $\nu(a)$ la multiplicité réduite de l'unique sommet de \ill{}
\note{le cadre entourant ces deux lignes est tracé à main levée ; la phrase suivante est écrite à l'encre rouge}

h) \struck{Pour tt arête $a$ de $K$, soit $\nu(a)$ le nb des pts simples de $L$ sur $a$ ($= \nu(s) + \nu(s')$, si $a$ est ordinaire d'extrémités $s$, $s'$) et} Soit $n$ l'ordre de $\Sigma$. Alors pour tt face $F$ de $K$, le nb des pts de $L$ sur $\partial F$, chacun compté avec sa multiplicité réduite, est $2n$.

\emph{NB} si $F$ ordinaire, \ill{} \ill{} :
\[
\sum_{s \text{ sommet de } F} \bigl(2\nu(s) + \varepsilon(s)\bigr) + \sum_{a \text{ arête de } F} \nu(a)
\]
\note{la formule est entourée d'un trait qui la sépare des figures ; la ligne qui l'introduit se lit mal}

\page{44}
(h) 1) Les faces de $K$ sont indexées par les p.r.\ pour $\underline{\Sigma}$ (à l'exclusion des p.r.\ \struck{exceptionn.} supercritiques),

2) Les arêtes de $K$ sont indexées par les \struck{paires} \add{couples} de p.r.\ \struck{\ill{}} $(D, D')$ telles que $D \xrightarrow{\sigma} D'$ (spécialisation immédiate), \struck{$\underline{\Sigma}$} \add{\ill{}} les flèches du \uncertain{$1$-graphe} \struck{entre positions} \struck{\ill{} \ill{} \ill{} \ill{}} : $F_D$ et $F_{D'}$ \struck{\ill{} \ill{} voisines d'une} \struck{\ill{} entendu que} $D \xrightarrow{\sigma} D'$ entre positions relatives voisines d'une $D_i$ \struck{\ill{}} de $\Sigma$ (\ill{} \ill{} $\Sigma$, $\Sigma'$ \ill{} \ill{} \ill{}, \ill{} \ill{} \ill{}). Si \uncertain{$a$} est \ill{} \ill{} : $F_D$ et $F_{D'}$, et elle est ordinaire ssi $\sigma$ est une spécialisation. Dans ce cas, le bord privilégié de $a$ \ill{} de d) est celui qui correspond à la face $F_D$.
\note{« (h) » est entouré, et marqué dans la marge de trois traits horizontaux. Le passage biffé entre « immédiate), » et « $D \xrightarrow{\sigma} D'$ entre positions » tient sur trois lignes et se lit mal}

3) Les sommets \add{(ordinaires)} \struck{\ill{}} de $K$ sont indexés par les \add{quadruples} \struck{triples} $(D, t', t'', \beta)$, où $D$ est une position relative \add{(non supercritique)}, $t$, $t'$ deux triangles de spécialisation pour $D$ \add{(\ill{} \ill{} par \ill{} sommet\ldots)} \add{(réels)}, et $\beta$ un \add{arc} \struck{\ill{}} de $\widetilde{D}$ (\ill{} des \ill{}), compris entre les arcs de spécialisation \struck{\ill{}} correspondant à $t$, $t'$. \ill{} \struck{Aux} \ill{} paires $(D, \beta)$, où $D$ est une p.r., et $\beta$ un arc de $\widetilde{D}$ \struck{\ill{}} \add{\uncertain{adjacent} des deux côtés à un arc de} spécialisation ordinaire). \struck{Le sommet $s$} \struck{cette donnée est \ill{} \ill{} \ill{} « sommet » $s$ de $D$} (ce qui implique que $D$ est stable). \add{\ill{} \ill{} \ill{} \ill{} \ill{}} \struck{\ill{} \ill{}} Le sommet $s$ correspondant a $4$ positions relatives $(D, D', D'', D''')$ : $D'$, $D''$ \add{$D'''$} \struck{\ill{}} sont distinctes \add{par spécialisation} de $D$ en \ill{} \ill{} $D \cap t'$, \struck{\ill{}} \ill{} $t''$, $D'''$ \struck{\ill{}} \ill{} les deux à la fois. \ill{} \ill{} \add{\uncertain{consécutives}}
\note{les lettres des triangles varient sur la page ($t'$, $t''$ dans le quadruple, $t$, $t'$ ensuite) ; transcrites telles quelles. La lettre $\beta$ est écrite chaque fois sur une autre lettre}

4) Les sommets exceptionnels sont indexés par les repères \ill{} $\coprod D$, \struck{i.e.\ \ill{} \ill{}} $=$ les \struck{\ill{}} \ill{} qui \ill{} \add{\ill{} arêtes orientées de} \struck{paires} \struck{triples} \struck{$(s, a)$ \ill{} \ill{} type \ill{}} $\underline{\Sigma}$, \struck{À une \ill{}} i.e.\ les drapeaux $(s, a)$ du type $(0, 1)$. \struck{À un \ill{} drapeau} Les \add{arêtes} \struck{\ill{}} issues de ce sommet $s$ sont les positions relatives voisines de $D_i$ \struck{\ill{} \ill{}} qui passent par $s$, et qui \struck{\ill{}} coupent $D_i$ en $s$ \ill{} $a$. Les quatre arêtes incidentes en conséquence.

\emph{NB} Ces indications suffisent à construire le tout cellulaire $(\mathcal{X}, K)$, qui est une géométrie d'incidence.

\emph{Cor.} L'ens.\ des comp.\ conn.\ de $K_0$ est indexé par $I$. \struck{Soit} \struck{$K(i)$} \add{\uncertain{la}} comp.\ conn.\ $K(i)$ ($i \in I$) est combinatoirement isomorphe à \struck{\ill{}} $\widetilde{D}_i$, l'image \ill{}

\page{45}
inverse de $D_i$ dans $\widetilde{X}$ (\ill{} \ill{} \ill{} au revêtement des rives de $D_i$, \ill{} \ill{} par choix d'orientation de $D_i$)

(i) On a une application
\[
\mathcal{X} \xrightarrow{\ \varphi\ } X
\]
(application structurale \add{\uncertain{contraction}}) telle que

1) La restriction de $\varphi$ à toute face de $(\mathcal{X}, K)$ est \ill{} \ill{}

\struck{2) \ill{} \ill{} \ill{} \ill{} $\varphi$ \ill{} \ill{} \ill{} $\mathcal{X}$, $K$}

2) Au voisinage des \struck{pts} \add{\ill{}} arêtes de $K$, \struck{\ill{} \ill{} \ill{}} $\varphi$ a la structure d'un « folding » d'arête.

3) Au voisinage des \struck{\ill{}} sommets de $K$, $\varphi$ a la structure d'un « double folding » \ill{} : $\xi$ et les deux branches de $K$ passant par $s$\ldots
\marginal{\ill{} \ill{} \ill{} \ill{} \ill{} ordinaires \ill{} \ill{}}
\note{« folding » et « double folding » sont écrits entre guillemets, en anglais. La note marginale, écrite en biais dans la marge de gauche à la hauteur de 2) et 3), se lit mal ; une ligne biffée la précède}

4) \struck{$L = \ldots$ \ill{} \ill{} \ill{} $\varphi^{-1}(\bigcup D_i) \cap K = S(L) \cap$} \quad $L = \varphi^{-1}\bigl(\underbrace{\textstyle\bigcup D_i}_{X_1}\bigr)$ \quad \struck{($\Rightarrow \varphi^{-1}(X_1) \cap K$)} \quad $K \cap L = (\varphi|K)^{-1}(X_1)$, \struck{\ill{}}
\note{« $X_1$ » est écrit sous l'accolade, sur un mot biffé}

\emph{Cor.}\ Un pt \struck{\ill{}} $s \in K \cap L$ \ill{} $\varphi(s)$ est un sommet, ssi $s$ \struck{est un sommet} \struck{\ill{} \ill{} \ill{} critique de} i.e.\ est un \emph{point critique} de $K$ (l'un des \ill{} pts est un \ill{} 1-1 \ill{} l'une des arêtes ordinaires de $K$).

(j) \struck{\ill{} Soit} Soit $F$ une face de $(\mathcal{X}, K)$. Considérons les \struck{sommets} \add{pts} critiques $s$ de $K$ sur $\partial F$, tels que le bord marqué \struck{\ill{}} $s$ (cf.\ (d)) soit extérieur à $F$, et corrigeons le contour $\partial F$ par un petit renflement \struck{\ill{}} au voisinage de $s$ \add{du côté} privilégié de $s$. Soit \struck{$F^{*}$} $F'$ \struck{la} la \struck{face \ill{}} ainsi définie \add{($F' \supset F$)} \struck{\ill{}} \struck{Alors} Alors la position relative correspondante à $F$ est $\varphi(\partial' F)$, et \struck{\ill{}} la \uncertain{situation} affine correspondante est canoniquement isomorphe à \struck{\ill{}} \struck{$F'$ \ill{}} $(F', F' \cap L)$.
\note{« (j) » est entouré et écrit sur un premier « (i) Soit » biffé}
\marginal{\emph{NB} (voir j) pour \ill{} pt critique $s$ de $F$ \ill{} \ill{}, i.e.\ \ill{} \ill{} privilégié \ill{} \ill{} \ill{} branches \ill{} \ill{} toutes \ill{} \ill{} \ill{} au bord \ill{} $F$ (dans le cas contraire \ill{} \ill{} \ill{} distinctes\ldots) \ill{} bord de $F$ \ill{} \ill{} $F$ et $F'$ \ill{} dessous}
\marginal{\ill{} \ill{} il faut \ill{} \ill{} « couloirs » \ill{} compris \ill{} de $K_0$}
\note{deux notes écrites en biais dans l'entre-colonne, la première reliée par une flèche à une figure. Figures : deux droites vertes et deux traits orange se croisant, marqués $s$, un trait bleu gris au point de croisement ; un octogone vert en deux moitiés, $F$ et $F'$, traversé d'un faisceau orange en forme de lentille resserrée au milieu en $s$ (un petit trait bleu), les pointes marquées $t$, $t'$, $t''$, $t'''$}

k) \struck{\ill{}} Une arête exceptionnelle \emph{stable} ou instable (ou semi-stable) suivant que les deux \struck{positions relatives} \add{positions} incidentes correspondantes sont des p.r.\ stables (resp.\ instables) (condition \ill{} resp.\ \ill{}). Alors pour \struck{\ill{}} \ill{} sommet exceptionnel, le bord privilégié de la branche ord.\ de $K$ qui y passe est du côté de l'arête exc.\ instable \add{(\ill{} \ill{})}

\page{46}
l) Loi de symétrie autour d'une arête exceptionnelle (\ill{} \ill{} \ill{} \ill{}), et structure locale autour d'une composante exceptionnelle de $K$
\note{figure sur toute la largeur de la colonne : une longue droite horizontale (le cercle exceptionnel), alternativement vert foncé et brun, coupée par des droites verticales vert clair marquées de petits traits bleus ; entre elles, des faisceaux de pseudo-droites orange qui se croisent au-dessus et au-dessous de l'axe en fuseaux et en lentilles, l'un hachuré au crayon noir ; au-dessus, deux petites étoiles orange sur des traits verts et bleus. Des traits au crayon relient la figure aux notes voisines}

figure canonique, \struck{il y a} \add{\ill{}} \ill{} j) \ill{}

\emph{NB} l'indice de chaque arête exceptionnelle stable est $= n - (\nu_{s_1} - 1) - (\nu_t - 1)$ \struck{\ill{}}, \add{et} les instables $= n - \nu$, où $\nu$ est la multiplicité du sommet --- dans la figure ci-contre on constate ! On voit que l'indice des arêtes exceptionnelles résulte des autres données numériques
\note{les signes $=$ sont appuyés, peut-être des $\equiv$ ; l'indice du premier $\nu$ est \uncertain{peu net}. Les dernières lignes de la note sont entourées d'une accolade, qu'une flèche relie à la figure}

m) Indexation des sections de $L$ sur les faces de $(\mathcal{X}, K)$ \struck{à l'aide des} par $\widetilde{I}$, à \ill{} près d'une indexation ; l'indexation de la face adjacente

n) Reconstitution de \struck{$\Sigma$} $\underline{\Sigma}$ (via $\widetilde{I}$\ldots) à l'aide de $K$, et des conditions numériques \add{$\varepsilon(s) \in \lbrace 0, 1 \rbrace$ et $\nu(s) \in \mathbb{N}$} \ill{} : $K$ (indices des sommets \add{ord.}\ de $K$), $\nu(a)$ des arêtes \add{\ill{}} de $K$ \struck{\ill{}} (distinguer le cas des arêtes \add{\ill{}} ordinaires et des arêtes exceptionnelles.)
\note{dans l'ajout sur $\varepsilon(s)$, le premier élément de l'ensemble est surchargé, lu $0$ --- lecture \uncertain{probable}}

\emph{NB} Il y a aussi évidemment une loi de symétrie autour d'une arête ordinaire, pour ce qui est de la configuration induite par $L$ sur les $2$ faces adjacentes, au voisinage de l'arête en question. Peut-être un énoncé \ill{} : considérons le couple $(\mathcal{X}, K \cup L)$, et « le voisinage tubulaire \ill{} » d'« \ill{} d'arêtes » \add{\ill{} \ill{}} de $K$ \struck{des arêtes} $K \cup L$ dans \ill{} arêtes. \struck{\ill{}} Alors pour toute arête $a$ de $K$, posons
\[
T_a \overset{\mathrm{def}}{=} (F' \cup F'') \cap T,
\]
où $F'$, $F''$ sont les deux faces de $(\mathcal{X}, K)$ incidentes à $a$, on a une involution $\sigma_a$ de $T_a$, \ill{} $a$ comme \ill{} \ill{} de pts fixes, et respectant la structure cellulaire induite par $(\mathcal{X}, K \cup L)$. De plus, \struck{si $a$ et $b$} \ill{} sont incidentes au \ill{} sommet $s$, alors $\sigma_a | T_a \cap T_b$ et $\sigma_b | T_a \cap T_b$ commutent dans un sens évident
\note{figure au milieu de la colonne : deux droites vertes (horizontale et verticale) traversées de pseudo-droites orange qui s'entrecroisent, un point de concours en pointillé orange}

\page{47}
\note{« 7.1.1984 » est écrit en haut à gauche : date de sa main (7 janvier 1984)}

\emph{Structure de $K$}. En rajoutant des sommets d'ordre $2$, soit $S_K$ l'ens.\ des sommets. On désigne par $K$ le graphe top.\ \ill{} ses nouveaux sommets (les anciens sont ceux d'\struck{\ill{}} ordre $4$). On a
\[
S_K = \underbrace{S'_K}_{\text{sommets rouges}} \amalg \underbrace{S''_K}_{\text{sommets noirs}}
\]
\note{sous $S''_K$, « noirs » est écrit au-dessus d'un « rouges » biffé}
\marginal{\emph{NB} \struck{\ill{}} $S'_K$ contient \ill{} sommets d'ordre \ill{} \ill{} des sommets \ill{} $S''_K$ contient des sommets \ill{} pas nécessairement d'ordre \ill{}}
\note{note écrite à droite de la formule, dans la colonne de gauche, en partie coupée par le bord de la zone lue}

\ill{} :

a) Pour toute face $F$, $S'_K \cap \partial F$ et $S''_K \cap \partial F$ \ill{} des subdivisions de $\partial F$ dont la \ill{} \ill{} des \ill{}, d'ordre $2n$
\struck{\emph{NB}} --- Deux sommets consécutifs sur $\partial F$ ne sont pas d'ordre $4$, et plus précisément, \struck{\ill{} des} \ill{} un arc de $\partial F$ joignant deux sommets d'ordre $4$, il y a toujours au moins un sommet rouge d'ordre $2$ (\ill{} \ill{} \ill{} si $\Sigma$ n'est pas \struck{\ill{}} du type $\mathrm{II}_p$ \ill{} \ill{} \ill{} \ill{} arête exceptionnelle) et au moins un sommet noir \ill{} \ill{} \ill{} B) \ill{} \ill{} \ill{} \ill{} \ill{} \ill{} \ill{}\ldots)
\note{« rouge », « noir » : les sommets sont effectivement marqués de ces couleurs sur les figures du dossier. La lettre « B) » est \uncertain{peu nette}}

b) \struck{\ill{}} On a une involution \ill{} \ill{} \ill{} \ill{} \ill{} sur le déploiement $K'$ de $K$. Ainsi, chaque arête $a$ d'une face $F$ est munie d'un signe $\varepsilon_F(a)$, qui est $\pm 1$ si le \ill{} est \ill{} \ill{} de $F$, $-1$ sinon. On a les règles
\[
\left\lbrace
\begin{array}{ll}
\varepsilon_{F'}(a') = -\varepsilon_F(a) & \text{si } F, F' \text{ sont les deux faces incidentes à } a \\
\varepsilon_F(a) = \varepsilon_F(b) & \text{si } a \text{ et } b \text{ sont adjacentes sur } F, \text{ le sommet commun étant d'ordre } 2 \\
\varepsilon_{a_1}(F_1) = \varepsilon_{a_2}(F_2) & \text{si } F_1 \text{ et } F_2 \text{ sont adjacentes le long d'une arête}
\end{array}
\right.
\]
mitoyenne $b$, d'extrémité $s$, d'ordre $4$, et $a_1$ et $a_2$ sont incidentes à $F_1$ et $F_2$ resp.\ \ill{}, $\neq b$, et incidentes à $s$.
\note{dans les deux premières règles, les lettres entre parenthèses sont surchargées ($a'$ sur une autre lettre, puis deux lettres noircies lues $a$, $b$ --- lecture \uncertain{probable}). Les conditions des règles sont écrites sur plusieurs lignes dans la colonne de droite ; nous les resserrons dans le tableau, la fin de la troisième étant rendue à la suite. Figure à côté de la troisième règle : une croix verte, un point noir au centre marqué $s$, une branche marquée $b$}

Il résulte du fait que toutes les faces sont de $=$ ordre \ill{}, que \struck{\ill{}} pour toute arête $a$, si $F$, $F'$ sont les deux faces incidentes, on a une \emph{bijection de symétrie}
\[
\sigma_a^F \text{ ou } \sigma_a : \operatorname{Pol}(F) \xrightarrow{\ \sim\ } \operatorname{Pol}(F')
\]
\marginal{\emph{NB} On a $\sigma_a^F \sigma_a^{F'} = \mathrm{id}$}
caractérisée par la condition qu'elle laisse \struck{\ill{}} \add{fixes} les deux extrémités de $a$. Il laisse donc \struck{\ill{}} \add{fixes} \struck{\ill{}} tous les sommets et arêtes qui sont sur la « grande arête » \add{\ill{}} contenant $a$ (\ill{} par les sommets d'ordre $4$ les plus proches) (mais \struck{\ill{} \ill{}} les prochains sommets \ill{} du \ill{} \ill{} \ill{} \ill{} \ill{} ne sont pas fixés par $\sigma_a$). L'involution $\sigma_a$ transforme sommets rouges (resp.\ noirs) en sommets noirs
\note{« bijection de symétrie » est souligné. Une note marginale écrite en biais sur le bord gauche de cette colonne, en partie coupée, n'est pas lue}

\page{48}
\note{« 2.1.84 » est écrit en haut à gauche : date de sa main, le premier chiffre est le même qu'à la page 50, où la date se lit nettement « 2.1.84. » (la page 47 porte « 7.1.1984 »). Deux longs traits obliques au crayon (l'un brun) traversent la colonne de haut en bas, sans masquer le texte ; la moitié droite de la feuille est blanche}

On place sur $K$ des nouveaux sommets (d'ordre $2$), de la façon suivante, \struck{Pour \ill{} \ill{} \ill{} \ill{} (\ill{} $F$) \ill{} \ill{} les sommets} arête par arête. \ill{}

a) On rajoute les pts de $K \cap L$ sur $\partial F$

b) Pour toute arête des nouveaux complexes topologiques \struck{\ill{}} $(\partial F, S \cap \partial F \cup (\partial F \cap L))$, \struck{\ill{}} on distingue
\note{« topologique distingue » est encadré, le cadre remontant jusqu'à a)}

\quad a) \ill{} rajoutant les pts de $\mathring{a} \cap L$

\quad --- soit $\underline{a}' = (a, a \cap S(K) \cup a \cap L)$
\note{figure à gauche : une ligne brisée verte (le bord $\partial F$), coupée de petits traits orange ; deux croix orange, l'une marquée $s$, l'autre au coin, le bord au-delà hachuré}

le graphe \add{topologique} \ill{} \ill{}, \ill{} $\alpha$ \ill{} \ill{} ce graphe. (\emph{NB} \struck{\ill{}} il y a au moins un \struck{\ill{} \ill{} \ill{}} élément de $a \cap L$ (« \struck{\ill{}} \add{sommet} rouge ») \ill{} a distinct des extrémités, donc chaque arête $\alpha$ de $\underline{a}'$ \struck{\ill{}} a au moins une extrémité rouge, i.e.\ qui \ill{} pas une \struck{\ill{}} \add{extrémité} de $\underline{a}$, et il y a \ill{} \ill{} \ill{} qui n'en ont \ill{} \ill{} \ill{} (\ill{} \ill{} \ill{})). a) Sur ces dernières, \add{(sauf si l'extrémité \ill{} est un sommet \ill{})}, on \ill{} \ill{} \ill{} \ill{} sommet. \add{b) Pour les} \struck{Sur \ill{}} autres arêtes, on \struck{\ill{}} rajoute des sommets (oranges) de la façon suivante : b) si \struck{\ill{}} les \add{deux} extrémités sont non critiques \ill{} \ill{} \ill{} des sommets de $L$, on rajoute juste un sommet (le milieu) c) \struck{Si \ill{}} Pour les \ill{} autres arêtes (cas d'une arête $\underline{a}$ ordinaire), de part et d'autre du pt critique $s$, si $\nu(s)$ est \ill{} ordre, on dispose $\frac{\nu(s)}{2}$ pts rouges (si
\note{la page s'arrête sur « (si » ; la numérotation a), b), c) est reprise à l'intérieur de b), telle quelle}

\page{49}
[j'aimerais comprendre la relation entre les $\sigma_a^F$, et les signes $\varepsilon_F(a)$.]
\note{le crochet ouvrant n'est marqué que par un grand trait vertical à gauche de la ligne}

Mais il faut encore introduire une autre structure

c) Ensemble $C_K \subset S_K$ des « pts critiques » \add{(ils sont d'ordre $2$, et tels que)} avec la condition que \add{sur} chaque \struck{\ill{}} grande arête, il y a \struck{exactement} un \struck{tel} élément de $C_K$ \struck{(quand il y en a pas, l'arête est exceptionnelle)}. De plus, chaque pt critique $s$ est muni d'un poids $\nu(s) \geq \struck{1}$, \add{\uncertain{premier} pair ou impair, suivant que $s$ est noir ou orange} (le poids d'un \ill{} de $S_K$ est par définition $\nu(s)$ si \ill{} critique, $1$ si non critique \struck{\ill{} \ill{}} et dans $S'_K$, $0$ s'il est non critique et dans $S''_K$)
\note{« ex. », « exactement » est écrit au-dessus de « \struck{\ill{}} » biffé ; « (quand \ldots{} exceptionnelle) » est biffé d'un trait. Le signe $\geq$ est suivi d'un $1$ surchargé, \uncertain{peut-être} un $0$ ; l'ajout interlinéaire sur la parité est souligné et se lit mal}
\marginal{définir les « \ill{} critiques » \ill{} « \ill{} de transition » \ill{} \ill{} \ill{} \ill{}, d'\ill{} \ill{} $\nu(s)$ sommets \ill{} $2\nu(s)$ (\ill{} \ill{} \ill{} \ill{} rouges)}
\marginal{\emph{NB} les \ill{} de transition (\ill{}) \ill{} \ill{} \ill{} \ill{} \ill{}\ldots}
\note{plusieurs lignes écrites en biais dans la moitié droite de la feuille, dont une série de notes à peine lisibles (« \ill{} critiques \ill{} \ill{} \ill{} \ill{} \ill{} \ill{} », « les sommets \ill{} \ill{} \ill{} des \ill{} critiques \ill{} », « relation entre $C_K$ et les signes $\varepsilon_F(a)$ ») ; nous n'en donnons que ce qui se lit}

En termes des sommets critiques, on a ceci :

1) Soit $s$ un sommet critique, incident à la face $F$, soit $s'$ son antipodique dans $\partial F$. \struck{Alors si $s'$} Pour qu'on ait \add{$\varepsilon_F(a) = +1$}, il faut et \ill{} \struck{\ill{} \ill{} critique \ill{} \ill{} le \ill{} $\varepsilon_a$} que le sommet $s'$ antipodique de $s$ (\add{\ill{} $\partial F$}) soit également critique (il a alors le $=$ indice)

2) Pour toute face $F$, il y a au moins $3$ pts critiques sur $\partial F$ (et $=$ au moins $4$,

3) \struck{Si toutes les pts} il y a au moins un pt critique\supplied{s} de signe $+1$. \struck{\ill{}} \add{\ill{} \ill{}}, sur chacun
\note{les rubriques 1), 2), 3) sont réunies par une grande accolade à droite. La parenthèse de 2) n'est pas refermée ; la page s'arrête sur « sur chacun »}

\page{50}
\note{« 2.1.84. » est écrit en haut à gauche : date de sa main (2 janvier 1984). Comme à la page 48, deux longs traits obliques au crayon traversent la colonne sans masquer le texte ; la page reprend, en la modifiant, la rédaction de la page 48}

On place sur \struck{chaque} $K$ deux ensembles de pts, $S'_K$, $S''_K$ \struck{\ill{}}, \add{(sommets rouges, sommets oranges)} de telle façon que

\struck{a)} $\forall F$ face de $K$, $\partial F \cap S'$ est un polygone \add{$\alpha$} top.\ d'ordre $2n$, et sur chaque arête \struck{$\alpha$} de celui-ci il y a un pt de $S''$ et un seul, qui est son milieu si $\alpha$ est contenu \ill{}
\note{cette condition est entourée d'un cadre, et reliée par une accolade aux conditions a) et b) qui suivent}

\struck{a) $S'_K \ldots$ \ill{} \ill{} les \ill{} sommets}

\struck{a) $S''_K \cap S(K)$, $S'_K \supset K \cap L \smallsetminus K \cap S(L)$}

a) $S'_K \cap L = K \cap L \smallsetminus$ (ensemble des pts critiques de $K$ \uncertain{pourvus} d'ordre \uncertain{pair})

\struck{b) $S'_K \cap S(K) = S(K) \cap L$}

b) \struck{Tout} Pour toute arête (\ill{}) $\alpha$ de $K$, on a $\alpha \cap S'_K \neq \emptyset$.
\note{à la suite de b), séparé par une flèche : « des \struck{\ill{}} arêtes de $K$, et l'unique sommet de $K$ sur $\alpha$ dans les \ill{} » ; figures : un segment vert barré de rouge et d'orange, biffé ; un segment vert en pointillé portant en son milieu une croix rouge ; un coin vert, un point rouge et un trait orange}

\struck{b) \emph{NB} a) et b) impliquent que $S'_K \cap S(K) = S(K) \cap L$, i.e.\ un \ill{} \ill{} de $K$ est dans $S'_K$ ssi il est dans $L$}

c) Pour un pt critique $s$ de $K$, sur l'arête \struck{$S'_K \smallsetminus S_K \cap L$ est réunion de segments, un pour chaque} \struck{pt \ill{}} \struck{ordinaire} $\alpha$ \struck{\ill{} qui} désignant par $t'$, $t''$ les pts les plus proches de $s$, sur l'arête $a$ de $K$ \ill{} $s$, parmi les pts de $(a \cap L) \cap \partial a$ \add{l'intersection $[t', t''] \cap S'_K$ est de cardinal égal} \struck{on dispose \ill{} le segment $[t', t'']$} $\nu(s)$, \struck{et il y a \ill{}}
\note{la phrase se termine sur la virgule après $\nu(s)$}

\note{figures dans la moitié droite de la feuille : en haut, un cercle noir et gris traversé de courbes qui se croisent, deux points verts ; deux hexagones et un pentagone rouges, chacun traversé d'un segment vert joignant un sommet (point rouge) à un côté ; un hexagone rouge partagé par un segment vert, hachuré au crayon ; au bas, un réseau de droites rouges qui se croisent en losanges, parcouru d'une ligne brisée verte, et un faisceau rouge coupé par des segments verts en pointillé}

\page{51}
\struck{\emph{NB}} des deux arcs \add{ouverts} \struck{déterminés} par $s$, $t$ sur $\partial F$, il y a au moins un pt critique

d) \emph{Relation entre} sommets critiques \struck{\ill{}} sur $F$ et sur $F' = \sigma_a(F)$ \add{(\ill{} $\partial F$. OPS que $\varepsilon_F(a) = \pm 1$)}, $a$ une \emph{grande arête}, sur $\partial F$ \add{(\ill{} les rôles de $F$ et $F'$)}. On a (\ill{} : \ill{} \add{$s$, l'antipodique de $s$, qui est critique}) \struck{\ill{} \ill{} \ill{} critique \ill{} \ill{}} \ill{} $u$, $v$ \ill{} ceci : Soit \struck{$s$ le pt critique sur $a$} \ill{} les deux pts critiques les plus voisins de $s$, \add{de part et d'autre} \struck{sur $\partial F$} l'arc correspondant \ill{} \ill{} de $s$, \ill{} $\beta = \widehat{u^{*} s\, v^{*}}$ \struck{\ill{}}. Considérons \struck{\ill{}} l'autre \add{\ill{}} d'extrémités $u'$, $v'$, contenant $s$. \struck{Soient} (\emph{NB}, \ill{} \ill{} extrémités $u'$ et $v'$, sont \add{\uncertain{connexes}} distinctes déterminées par $a$, $u$, $v$, par \ill{} \ill{}, \ill{} $u \neq v$, \ill{} $s \in \beta$, \ill{} $s$, $t$ sur $\partial F$, \struck{\ill{}} et $u \neq v$, \ill{} \ill{} \ill{} $t'$, $u \sqcup v'$ \ill{} façon précise $\alpha = \widehat{u s v}$). \struck{\ill{}} \ill{} transformés de $s$, $t$, $u$, $v$ par $\sigma_a$. Alors

a) \struck{\ill{}} \struck{\ill{} induit une bijection \ill{} \ill{} \ill{} condition, \ill{}} \emph{bijection entre} $C_K \cap \alpha$ et $C_K \cap \alpha'$, \ill{} $x$ et $x'$ se correspondant \ill{} $\varepsilon_F(x) = \varepsilon_F(x')$, sauf si $x = s$, donc $x' = s$, \ill{} \ill{} $\varepsilon_{F'}(x) = -\varepsilon_F(x)$ (\ill{} \ill{} \ill{}\ldots)
\note{« Relation » (avec « entre ») et « grande arête » sont soulignés ; « $C_K \cap \alpha$ et $C_K \cap \alpha'$ » est souligné d'un trait. Le passage entre « Considérons » et « Alors » est chargé d'ajouts interlinéaires et de ratures ; les lettres $u^{*}$, $v^{*}$ portent un astérisque \uncertain{peu net}}

b) On a $t' \notin C_K$, \struck{\ill{}} \ill{} \ill{} \ill{} $F$ \ill{} $F'$ \ill{} $N \ldots N'$, \ill{} \ill{} $\operatorname{card}(\mathring{\beta}' \cap C_K) \in \lbrace 0, 1, 2 \rbrace$, $N' \in \lbrace N-1, N, N+1 \rbrace$ \struck{\ill{}} \struck{[\ill{} si l'indice de \ill{}]} si $x \in \mathring{\beta}' \cap C_K$, on a $\varepsilon_F(x) = -1$
\note{b) est surchargé : une formule biffée précède « $\in \lbrace 0, 1, 2 \rbrace$ » ; « $N' \in \lbrace N-1, N, N+1 \rbrace$ » est écrit au-dessous ; le « $\mathring{\beta}'$ » est lu d'après la ligne suivante}

\emph{Cor.} Considérons les antipodiques $u^{*}$, $v^{*}$ de $u$, $v$ sur $\partial F$, et l'arc $\beta^{*} = \widehat{u^{*} s\, v^{*}}$, \struck{Alors} alors \ill{} \ill{} \ill{} il n'y a pas de pts critiques de signe $\pm 1$
\note{figure dans la marge gauche de la colonne : deux cercles verts tangents en un point marqué $s$ ; sur le cercle du haut, les points $u$, $t$, $v$ (points noirs et verts), l'intérieur griffonné au crayon ; sur celui du bas, $u'$, $t'$, $v'$}

\emph{En résumé}, soit $x \in S_K \cap \partial F$, \struck{\ill{} un sommet de $s$}, $x' = \sigma_a(F)$, \add{alors} $x'$ est critique \ill{} \ill{} \ill{} \ill{} \ill{}, et $x$ et $x'$ \ill{} \ill{} \ill{}, \ill{} \ill{} \ill{} les cas suivants
\note{« En résumé » est souligné d'un trait ondulé ; « $x' = \sigma_a(F)$ » est écrit tel quel, pour $x' = \sigma_a(x)$ vraisemblablement}

a) $x = s$, auquel cas $\varepsilon_{F'}(x) = -\varepsilon_F(x')$ \ill{} \ill{} critiques

b) $x = t$, auquel cas \ill{} $x'$ \ill{} critiques d'indice $\pm 1$, \ill{}

(Ceci permet de distinguer \struck{\ill{}} \add{\uncertain{laquelle}} \add{« \uncertain{privilégiée} »} \struck{« \ill{} »} des deux faces $F$, $F'$ \ill{} \ill{} \ill{} \ill{}, en termes de la position de l'\ill{} critique \ill{} \ill{} \ill{} \ill{} antipodiques \ill{} \ill{} $\partial F$, $\partial F'$)

c) $x \in \widehat{u h v} = \mathring{\beta}$, \ill{} $x \neq t$, i.e.\ $x' \neq t'$, donc $x$ \ill{} critique, alors \ill{} \ill{} \add{\ill{}} \ill{} \ill{} (\ill{} \ill{} \ill{}) \ill{} \ill{} \ill{} \ill{} $x'$ sont critiques \ill{} \ill{} indice est $-1$.
\note{« $\widehat{u h v}$ » : la lettre du milieu est \uncertain{peu nette}, lue $h$}

\page{52}
\note{page de dessins, sans texte de sa main : sur la moitié inférieure droite, deux figures séparées par un grand arc rouge. Chacune montre des faisceaux de droites rouges concourant en un point, et une ligne brisée verte (le bord d'une face) qui serpente entre elles, portant des sommets marqués de points noirs ou orangés ; dans la figure de gauche, un segment vert foncé relie le point de concours à un sommet, et une corde verte en pointillé longe le bas ; dans celle de droite, une corde verte en pointillé, jalonnée de points rouges, passe par le point de concours (point rouge épais). Ce sont vraisemblablement des illustrations du bord $\partial F$ et de ses sommets critiques, pages 49 à 51}

\page{53}
\emph{Réciproquement}, on \ill{} utiliser la position des \ill{} critiques, pour déterminer les signes des \struck{\ill{}} $\varepsilon_{F'}(a')$, quand on connaît les signes \struck{\ill{}} $\varepsilon_F(a)$.
\note{« Réciproquement » est souligné ; « les signes » (seconde occurrence) est écrit au-dessus d'un mot biffé et souligné}

\emph{Proposition}. Considérons l'arc \struck{\ill{}} $V_F$ (« vert ») des sommets d'ordre $4$ \struck{\ill{}} incidents à $F$, et l'arc $\sigma_a^{F'}(V_{F'})$ (où $F'$ est la face incidente à $F$ le long de la grande arête $a$ \struck{de} $\partial F$). Soit \struck{\ill{}} les transformés \ill{} \ill{} \ill{} de $\partial F$ \struck{\ill{} \ill{}} \add{arêtes} considérons les (fermées) $\alpha$, sont $A_{F,a} = $ \struck{\ill{}} la réunion des grandes arêtes \add{(sur $F$)} qui sont telles que $\sigma_a(\alpha)$ soit une grande arête de $F' = \sigma_a(F)$ (\emph{NB} $a$ est une grande arête de $F$). Alors

a) $A$ est un arc sur $\partial F$, soit $B$ l'arc « complémentaire ». \struck{Alors} Considérons la subdivision de $B$ en grandes arêtes (sommets verts clairs) et en « grandes arêtes » \ill{} $F''$ (sommets verts foncés). Alors le nb d'arêtes $n_a(F)$ et $n_a(F')$ qui interviennent est \ill{} \ill{} \ill{} compris entre $3$ et $5$ --- les possibilités sont (compte tenu de la symétrie \ill{} $F$ et $F'$) $(3, 3)$, $(3, 4)$, $(3, 5)$ [\struck{Donc \ill{}} (\ill{} $(2, 2)$ \ill{})] --- et dans les \ill{} cas, c'est la face spéciale qui donne lieu à $3$ arêtes (ce qui suffit à la déterminer, quand on est dans le cas $1^{\circ}$ ou $3^{\circ}$). De plus, \ill{} \ill{}
\note{« Proposition » est souligné ; « Soit » est suivi d'un mot biffé et noirci, la construction de la phrase restant incertaine. Les couples, surchargés, sont lus $(3, 3)$, $(3, 4)$, $(3, 5)$ --- les seconds chiffres, \uncertain{peu nets}, étant lus d'après « compris entre $3$ et $5$ » ; « Donc » et « \ill{} » sont biffés dans le crochet}

b) Si $\alpha$ est une \add{petite} arête dans $F$ dans $A$, \struck{\ill{} \ill{}} et $\alpha \neq a$, on a $\varepsilon_{F'}(\alpha) = \varepsilon_{F}(\alpha)$ [\emph{NB} si $\alpha = a$, on a $\varepsilon_{F'}(a) = -\varepsilon_F(a)$]

c) On a les positions suivantes \ill{} pour la distribution des grandes arêtes sur $B$ (à identifier près de $F$ et $F'$) et les signes correspondants
\note{figures dans la moitié inférieure de la colonne de droite, numérotées $1^{\circ}$), $2^{\circ}$), $3^{\circ}$) : des arcs verts (clair et foncé) d'extrémités $u$, $v$, portant des sommets $x$, $y$ (vert clair) et $r$, $s$ (vert foncé), avec les signes « $+1$ » au sommet de l'arc, « $-1$ » entre $r$ et $s$, et des accolades marquées $\varepsilon_1$, $\varepsilon_2$ de part et d'autre. Une première version de $1^{\circ}$) et $2^{\circ}$), au crayon, est barrée de grands zigzags et reprise au-dessous : en $1^{\circ}$), un seul sommet foncé $r$, de signe $+1$, entre $x$ et $y$ ; en $2^{\circ}$), deux sommets foncés $r$, $s$ séparés par un arc de signe $-1$, sous le $+1$ ; en $3^{\circ}$), trois sommets foncés, le médian de signe $+1$, deux arcs de signe $-1$ de part et d'autre}

\page{54}
\emph{Corollaire}. Dans tous les cas, la distribution des signes $\varepsilon_{\cdot}(F)$ détermine celle des $\varepsilon_{\cdot}(F')$. \struck{C'est clair \ill{}} C'est clair d'abord dans les cas $1^{\circ}$ et $3^{\circ}$, où on sait laquelle des faces $F$, $F'$ est la spécialisée de l'autre, \struck{et par suite} \add{Mais} dans le cas $(2^{\circ})$ \add{« symétrique »}, \struck{\ill{} \ill{} on connaît les signes $\varepsilon_{\cdot}(F)$, on} On voit que \struck{les} \add{grandes} deux arêtes extrêmes correspondantes pour $F$, $F'$ ont $=$ signes (c'est vrai dans tous les cas, y compris dans les cas $1^{\circ}$ et $2^{\circ}$), et pour les arêtes qui \add{\ill{}} sont par deux en \ill{}, elles sont \ill{} $F$, $F'$ par \ill{} \ill{} de $=$ signes (\emph{NB} leur nb, suivant les cas, est $(1, 0)$, $(1, 1)$ ou $(1, 2)$) et les signes en question pour $F$ et pour $F'$ sont opposés.
\note{« Corollaire » est souligné ; « symétrique » est écrit au-dessus de la ligne, entre guillemets, et « Mais » au-dessus de « et par suite » biffé}

\emph{Corollaire}. Quand on connaît les signes $\varepsilon_F(a)$ pour \struck{une} \add{\ill{}} face $F$, on les connaît pour \emph{toutes} (grâce à la connexité de $\mathcal{X}$)
\note{« une » est biffé d'un double trait ; « toutes » est souligné}

J'ai supposé implicitement \struck{que} dans les vérifications de la proposition que ni $F$, ni $F'$ n'était une face exceptionnelle (correspondant à une position relative exceptionnelle, i.e.\ « supercritique »). Il faudrait que j'examine séparément \ill{} cas de telles faces. Pour une telle face, il n'y a pas \ill{} \ill{} grandes arêtes, et la distribution des signes est \ill{} \ill{}. Cela semble donner une façon simple \add{(i.e.\ les \ill{} \ill{})} de décrire les signes \add{(\ill{} \ill{})} \ill{} structure de \ill{} \ill{} de $(\mathcal{X}, K, S_K)$, sans avoir à utiliser les \ill{} critiques, et encore moins \struck{les} le \struck{\ill{}} graphe $L \subset \mathcal{X}$.
\note{la colonne de droite occupe les deux tiers supérieurs de la feuille ; le bas de la feuille est blanc}

\page{55}
\note{page de calculs au crayon, sans texte suivi ; nous donnons les formules dans l'ordre de la feuille, colonne de gauche puis colonne de droite}

\[
(\nu_{s_1} - 1) + (\nu_{s_2} - 1) + \cdots + (\nu_{s_n} - 1)
\]
\note{la dernière parenthèse porte en dessous, dans une accolade, « $\nu_{s_0}$ ». Figures : un disque au crayon, à la périphérie duquel sont accolés un pentagone et un hexagone orientés (petites flèches courbes), deux points $0$ et $1$ sur le bord ; à droite, un segment vertical portant les points $s_1$, $s_0$, $s_{n-1}$}

\[
\begin{array}{l}
l_s = n - \nu_s \\
l_a = n - 1 - (\nu_s - 1) - (\nu_t - 1) = n - (\nu_s + \nu_t - 1) = (n - \nu_s) - (\nu_t - 1) \\
l_s - l_a = \nu_t - 1 \\
\chi_{2n}(s) = \sum (\nu(s) - 1) = n - 1
\end{array}
\]
\note{dans $l_s$, un premier coefficient devant $n$ est biffé ; dans $l_a$, un $2$ \uncertain{biffé} devant $n$, et dans la ligne suivante un terme biffé ; après $l_s - l_a = \nu_t - 1$ suit un « $= 1$ » biffé. Sous la somme, en accolade : « $s$ sommet sur $D_i$ » et « multiplicité du sommet », puis « $\mathbb{Z}/n\mathbb{Z}$ »}

\[
\lambda(t) = \widetilde{\rho}_p(t) = (-1)^p (t - \chi_p) \qquad \text{ici } p = 2n \text{ donc } (-1)^p = +1
\]
\[
\lambda(t) = t - \chi_{2n} = t - (n - 1)
\]
\note{devant la seconde formule, un signe \uncertain{peu net} (lu $\beta$ ou $\beta$ biffé) et « il faut » biffé}

\emph{NB}
\[
\rho_p(t) = \sigma_{p-1} \cdots \sigma_0 (t) = (-1)^p (t - \chi_p)
\]
\note{une première valeur, « $(-1)^{p-1}(\chi_p - t) = (-1)^p t + (-1)^{p-1}\chi_p$ », est biffée sur deux lignes}

Prolongement de $\varphi_p : F_0 \ldots$ fait par \struck{\ill{}} \ill{}
\note{la ligne se poursuit par des mots biffés, et une petite figure (un segment à deux points) ; nous ne la lisons pas au-delà}

\note{figure au milieu de la feuille : une chaîne verticale de faces polygonales $F_0$, $F_1$, $F_2$, \ldots, $F_{p-1}$, $F_p$, $F_{n+1}$, accolées le long d'un segment, les sommets de recollement marqués $a_0$, $a_1$, \ldots, $a_p$, $a_{p+1}$, $s_0$, $s_1$, \ldots, $s_n$, $s_{n+1}$, avec des flèches courbes ; « $p = 2n$ » à côté}

\[
\begin{array}{ccc}
\operatorname{P}(F_n) & \xrightarrow{\ \sigma_{a_n}\ } & \operatorname{P}(F_{n+1}) \\
\varphi_n \ \wr & & \wr\ \varphi_{n+1} \\
\mathbb{Z}/N\mathbb{Z} & \xrightarrow{\ \sigma_n\ } & \mathbb{Z}/N\mathbb{Z}
\end{array}
\qquad s_{n+1} \longmapsto s_{n+1}
\]
\note{« $\mathbb{Z}/N\mathbb{Z}$ » : le $N$ est \uncertain{peu net}, le second terme de la ligne du bas est surchargé. « $s_{n+1} \mapsto s_{n+1}$ » est écrit au-dessus de la flèche $\sigma_{a_n}$}

$\sigma_n$ a) respecte la structure polygonale

b) renverse l'orientation

c) transforme $\varphi_n^{-1}(s_{n+1})$ en $\varphi_{n+1}^{-1}(s_{n+1}) = 0$
\note{sous $\varphi_n^{-1}(s_{n+1})$, une accolade : « $l_n$ ». Figure : un cercle au crayon en pointillé, marqué « donc »}

\[
\sigma_n(t) = l_n - t
\]
\[
\begin{array}{ll}
\sigma_0(t) = l_0 - t & \\
\sigma_1(t) = l_1 - t & \sigma_1\sigma_0(t) = l_1 - (l_0 - t) = l_1 - l_0 + t \\
\sigma_2(t) = l_2 - t & \sigma_2\sigma_1\sigma_0(t) = l_2 - l_1 + l_0 - t
\end{array}
\]
\[
\underbrace{(\sigma_n \sigma_{n-1} \cdots \sigma_0)}_{\rho_{n+1}}(t) = l_n - l_{n-1} \cdots + (-1)^n l_0 + (-1)^{n+1} t = (-1)^n (\chi_{n+1} - t) = (-1)^{n+1}(t - \chi_{n+1})
\]
\[
\text{où } \chi_{n+1} = \underbrace{l_0 - l_1 + \cdots + (-1)^n l_n}_{n+1 \text{ termes}}
\]
\note{au-dessus de $(\sigma_n \sigma_{n-1} \cdots \sigma_0)$, en accolade : « $n+1$ facteurs » ; le $\sigma_{n-1}$ est surchargé. Au bas de la feuille : « $\underbrace{\sigma_{p-1}\sigma_p \cdots \sigma_0}_{\rho_p}(t) = (-1)^{p-1}\ldots$ », inachevé}

\page{56}
\[
\nu(n-1) + \nu(n-3) = 2\nu n - 4\nu = 2\nu(n-2)
\]
\note{la formule est écrite en haut à droite de la feuille ; sous « $(n-1)$ » un petit « $2$ » en accolade ; « $- 4\nu$ » est surchargé, et deux valeurs intermédiaires sont biffées, dont « $= \nu(2n-3)$ » ; le résultat final est encadré. Au-dessous, « $2(n-2)$ », surchargé}
\note{page de dessins pour le reste : trois « écuelles » vertes (un côté droit et un arc, jalonnés de sommets verts et noirs) accompagnées des couples $(3, 3)$, $(3, 4)$, $(3, 5)$, écrits à l'envers --- ce sont vraisemblablement les trois cas $1^{\circ}$, $2^{\circ}$, $3^{\circ}$ de la page 53 ; deux autres écuelles, l'une marquée $-1$ de part et d'autre, suivie de « $\ldots$ », et de nouveau « $(3, 3)$ $(3, 4)$ --- » à l'envers ; un cercle vert et noir avec des arcs épais noirs et des marques rouges, annoté $1/2$, $1/1$, $1$ ; en haut à droite, un pentagone aux sommets numérotés $0$ à $3$ et un quadrilatère, avec « $1\,2\,3$ » ; au bas, trois grands cercles au crayon, avec des arcs épais noirs et des points verts, l'un marqué $+$ en haut}

\page{57}
\note{« 3.1.84 » est écrit en haut à gauche : date de sa main (3 janvier 1984), le dernier chiffre étant lu comme à la page 50 --- lecture \uncertain{probable}. Deux longs traits obliques, l'un au crayon, l'autre jaune, traversent chaque colonne ; ils ne masquent pas le texte, et rien n'indique s'ils l'annulent}

Transport parallèle le long d'un demi-tour, \struck{\ill{}} le long d'une position relative \struck{singulière \ill{}} \add{supersingulière}
\[
\boxed{\lambda(t) = t - (n-1)}
\]
\note{cette formule reprend celle de la page 55, $\lambda(t) = t - \chi_{2n} = t - (n-1)$}
\marginal{Tout le \ill{} \ill{} \ill{} !}
\note{écrit en biais entre les deux colonnes}

\emph{NB} Le pgcd de $n-1$ et de $2n$ est $1$ si $n$ pair, $2$ si $n$ impair. Donc on voit que l'image de
\[
\pi_1(\mathcal{X}) \longrightarrow \mathbb{D}_{2n}
\]
a une image qui contient $\mathbb{D}^{+}_{2n}$ \add{$\simeq \mathbb{Z}/n\mathbb{Z}$} si $n$ pair (donc cette image est d'indice $1$ ou $2$) \struck{ou bien elle} et elle contient $\mathbb{Z}/n\mathbb{Z}$ si $n$ impair (donc elle est \ill{} d'indice $1$, $2$ ou $4$)
\note{« $\mathbb{D}_{2n}$ » est écrit sur un $\mathbb{Z}$ biffé. L'ajout « $\simeq \mathbb{Z}/n\mathbb{Z}$ », au-dessus de la ligne, est rattaché par une accolade à $\mathbb{D}^{+}_{2n}$}

\struck{Cela \ill{} \ill{} \ill{} \ill{} le \ill{} \ill{} \ill{}} \add{\ill{} \ill{} \ill{} le déployé $\widetilde{\mathcal{X}}$ de $\mathcal{X}$} suivant les composantes supersingulières de $K$, on trouve dans l'image de
\[
\pi_1(\widetilde{\mathcal{X}}) \longrightarrow \mathbb{D}_{2n}
\]
les \ill{}
\[
t \longmapsto \lambda'(t) = t - 2(n-1).
\]
Or le pgcd de $2(n-1)$ et $2n$ est \ill{} $2$. Donc l'image contient \ill{} $\mathbb{Z}/n\mathbb{Z}$, d'indice $4$
\marginal{\ill{} \ill{} \ill{} \ill{} \ill{} \ill{} \ill{} \ill{} \ill{} $\operatorname{Pol}(D)$ \ill{} \ill{} \ill{}}
\note{note écrite en biais dans l'entre-colonnes, séparée du texte de droite par un trait vertical ; elle se lit mal}

(le quotient étant $\mathbb{D}_2 \simeq \mathbb{F}_2 \times \mathbb{F}_2$)

Ceci ne nous dit \add{\ill{}} \ill{} l'image de $\pi_1(\mathcal{X})$ \ill{} $\pi_1(\widetilde{\mathcal{X}})$ dans $\mathbb{D}_{2n}$, et en particulier sur son image dans $\pm 1$, \add{\ill{}} \ill{} \ill{}
\[
\mathbb{D}_{2n} \xrightarrow{\ \chi\ } \pm 1
\]
On \ill{} l'impression que pour avoir des informations, il faut faire presque un \ill{} autour d'un sommet de $\underline{\Sigma}$ et faire \ill{} un \ill{} et voir l'effet produit. On va faire un demi-tour, \ill{} partir d'une position supersingulière \ill{} disons, \struck{\ill{}} transport parallèle \ill{} \ill{} l'orientation de la position supersingulière \ill{} changée, et \struck{\ill{}} $\chi(\lambda)$ sera égal à $+1$ si le nb total de « pas » pour revenir au pt de départ est \emph{impair}, à $-1$ dans le cas contraire. Donc on aimerait pouvoir le faire en un nb \emph{pair} de pas !

\page{58}
\note{un long trait oblique au crayon traverse la feuille en diagonale, sans masquer le texte}
Par une petite modification de ce \uncertain{processus} on montrera alors que non seulement $\pi_1(\mathcal{X}) \to \pm 1$, mais aussi $\pi_1(\widetilde{\mathcal{X}}) \to \pm 1$ est surjectif ; je dirai dans le cas $n$ pair, cela implique que $\pi_1(\widetilde{\mathcal{X}}) \to \mathbb{D}_{2n}$ est surjectif (mais peut-être pas $\pi_1(\widetilde{\mathcal{X}}) \to \mathbb{D}_{2n}$ ?)
\note{le premier $\widetilde{\mathcal{X}}$ de la dernière ligne porte au-dessus de son tilde un petit signe \uncertain{peu net}, peut-être un second tilde ou un « $\approx$ » : les deux revêtements comparés ne se distinguent que par ce signe}

Calcul de la parité \struck{\ill{}} du nb d'opérations élémentaires pour passer de $(D_i^{s}, \omega)$ à $(D_i^{s}, -\omega)$. Le nb d'opérations est égal à :
\[
\begin{array}{l}
2(\text{nb de sommets : balayage}) \\
\quad + 2(\text{nb de droites par-dessus lesquelles balayage}) \\
\quad + 1 \ (\text{décollage}) + 1 \ (\text{atterrissage})
\end{array}
\]
merveilleux, c'est toujours pair !
\note{figure à gauche : un segment horizontal portant un point $s$, une petite flèche verticale au-dessus. En regard de la deuxième ligne, une accolade renvoie à une note : « \emph{NB} chaque droite par-dessus laquelle on balaye compte comme un dessous sommet \uncertain{pratiquement} \ill{} balayerait »}

\struck{En fait, par transport parallèle, on trouve (\ill{} \ill{} \ill{} p.r.} À vrai dire, le raisonnement est \uncertain{incomplet}, \struck{\ill{}} car ce que il faut regarder ici dans le revêtement de $\mathcal{X}$ qui \ill{} par des faces supersingulières (par des « bandes supersingulières ») \ill{} celles-ci sont d'ordre $2(n-1)$, non $2n$. Donc il vaut mieux \ill{} partir d'une position proche de $D_i$ passant par $s$ par un petit pivotement autour de $s$. Maintenant, toute droite \ill{} franchie compte comme une \emph{seule} opération, \struck{\ill{}} \ill{} droites \add{\ill{}} franchies \add{(\ill{} \ill{})} \struck{\ill{}} une fois et une seule. Donc le nb correct d'opérations est
$2(\text{nb des sommets } \ldots\ s) + \text{ordre de } s$
\note{dans la parenthèse, deux mots entre « sommets » et « $s$ » sont illisibles : \ill{} \ill{}}
donc l'image du circuit dans $\pm 1$ est $-1$ ou $+1$, suivant que l'ordre $\nu(s)$ de $s$ est pair, ou impair. Si donc il y a au moins un sommet simple, on voit que l'homomorphisme
\[
\pi_1(\widetilde{\mathcal{X}}) \longrightarrow \pm 1
\]
surjectif, et \uncertain{aussi} $\pi_1(\widetilde{\mathcal{X}}) \to \pm 1$ surj.
\note{« seule » est souligné ; « l'ordre » est écrit sur « $\nu$ » biffé. La dernière ligne, au bord de la feuille, se lit mal ; le premier des deux $\widetilde{\mathcal{X}}$ porte de nouveau le signe \uncertain{peu net} au-dessus du tilde}

\page{59}
\note{la colonne de gauche est barrée d'un long trait oblique, et sa moitié inférieure de quatre traits parallèles au crayon ; la colonne de droite d'un trait oblique. Le texte reste lisible sous ces traits, et nous le donnons}
Par contre, si on fait un tour complet, de façon à ramener : l'orientation de \ill{} de $D$, \ill{} (i.e.\ le carré des \ill{} demi-tours) on trouve \struck{\ill{} \ill{} \ill{}} \add{un transport parallèle} $\operatorname{Pol}(D) \to \operatorname{Pol}(D)$ qui conserve l'orientation. \ill{} ainsi que \struck{\ill{}} pour les « \ill{} demi-tours autour de $s$ », \struck{\ill{}} \ill{} le nb de pas est pair, alors le résultat du transport parallèle sur les polygones de \ill{} \ill{}, sur l'orientation, les $=$ résultat que le transport parallèle ordinaire de l'orientation. Cela signifie que l'orientation change par le long d'un circuit, si $\nu(s)$ pair. Si $\nu(s)$ impair, alors le résultat \add{sur} $\operatorname{Pol}(D)$, et via transport ordinaire, n'est \struck{\ill{}} \ill{} les $=$ \struck{\ill{}} \add{\uncertain{comme}} l'orientation est \emph{conservée} par transport parallèle, elle ne l'est pas par transport ordinaire --- donc en tous cas, cette dernière change
\note{« conservée » est souligné}

\ill{} \ill{} chaque fois qu'on a un circuit fermé sur $\mathcal{X}$ (circuit \ill{} \ill{} \ill{} \ill{}) tel que le transport par \ill{} change l'orientation de la position relative $D$. \struck{Alors l'orientation \ill{} \ill{} \ill{} \ill{} de $\mathcal{X}$ \ill{} \ill{}} \ill{} \ill{} ssi le nb des pas élémentaires est pair. Balayer au-dessus d'un pt compte comme $2$ balayages, au-dessus d'une p.r.\ supersingulière compte comme $1$.
\note{« $\mathcal{X}$ » de la première ligne est précédé d'une barre verticale ; « de la position » est écrit sur une lettre biffée ($\Delta$ \uncertain{peut-être})}

Voilà enfin une façon de voir que $\pi_1(\mathcal{X}) \to \mathbb{D}_{2n}$ est surjectif.
\note{un mot biffé et noirci est écrit au-dessus de $\mathcal{X}$, peut-être un tilde}

On part d'une p.r.\ \add{$D$} stable proche d'une p.r.\ supersingulière $D_0$,
\note{figure au bas de la colonne : deux droites orange verticales et une droite orange horizontale $D_0$ ; une droite verte $D$, parallèle à $D_0$ au-dessus d'elle, s'infléchit pour la rejoindre et la croiser près du point $s$ ; le point $a$ est marqué sur $D$, et la région entre $D$ et $D_0$ près de $s$ est hachurée au crayon}

\page{60}
\struck{d'où un \ill{}} \add{d'où un} côté privilégié \add{\ill{}} sur $D$ (celui qui rencontre $D_0$), donc par ailleurs une extrémité $s$, de façon à avoir un repère. \ill{} spécialisation $D$ en $D'$ vers le \struck{\ill{}} \add{sommet} $s$ le plus proche de $s$, \ill{} fais un demi-tour par $s$ et je reviens : $D$, \ill{} un nb pair d'opérations, qui change l'orientation de $D$, \emph{et transforme $s$ en $s$} --- donc c'est la symétrie p.r.\ : $s$. Si on fait
\note{un trait oblique au crayon barre le haut de la colonne. « et transforme $s$ en $s$ » est souligné ; le « $s$ » après « symétrie p.r.\ : » est lu d'après la ligne précédente}

4.1. Pour tout \struck{\ill{}} chemin combinatoire $c$ sur $\Gamma$ --- i.e.\ chemin formé avec des faces $F_D$ de $(\mathcal{X}, K)$, \struck{\ill{}} \ill{} la face $F = F_D$ à $F' = F_{D'}$, on trouve un iso.\ « de transport parallèle »
\[
\operatorname{Pol}(c) : \operatorname{Pol}(F) \longrightarrow \operatorname{Pol}(F'),
\]
qui du pt de vue de la géométrie intrinsèque de $(\mathcal{X}, K, S_K)$ s'interprète par composition \ill{} de « réflexions » le long des arêtes de réflexion qui interviennent au cours de $c$.
\note{« 4.1. » ouvre l'alinéa, sans que rien sur ces pages ne donne les numéros précédents ; le $c$ est ajouté au-dessus de « chemin »}

L'effet de $\operatorname{Pol}(c)$ sur les orientations
\[
\underline{\omega}_{\ast}(c) : \underline{\omega}(F) \simeq \underline{\omega}(D) \longrightarrow \underline{\omega}(F') \simeq \underline{\omega}(D')
\]
est aussi l'effet sur l'orientation des p.r.\ par transport parallèle --- donc est \ill{} \ill{} \ill{} la parité des nb de pas dans le chemin $c$ ! Si on a un circuit $c$, i.e.\ $F = F'$ \ill{} trouve
\[
\pi_1(\mathcal{X}, F_D) \longrightarrow \underset{\simeq \mathbb{D}_{2n}}{\operatorname{Aut}(\operatorname{Pol}(D))} \xrightarrow{\ \text{effet sur orientation}\ } \pm 1
\]
\note{l'indice de $\underline{\omega}_{\ast}(c)$ est \uncertain{peu net} ; le « $\ast$ » est lu d'après le $\chi_{\ast}$ qui suit. Un « $\ell$ » est biffé devant « $\ell = c$ »}
On désigne par $\chi_{\ast}$ le caractère composé
\[
\chi_{\ast} : \pi_1(\mathcal{X}, F_D) \longrightarrow \pm 1
\]
donc pour un chemin $\ell$, on a $\chi_{\ast}(\ell) = +1$ ou $-1$ suivant que par transport parallèle l'orientation de $D$ \ill{} ou non changé. \struck{\ill{}}

On a deux autres caractères $\chi_s$ et $\chi_a$ (\struck{\ill{}} \add{orientation} « sommets » et « arêtes ») \ill{} \ill{} : chaque face $F_D$ \ill{} deux \ill{} \ill{} \ill{} d'éléments
\[
\begin{array}{l}
\varepsilon_s(F) = \varepsilon_s(D) = \text{un des deux } n\text{-gones inscrits dans } \operatorname{Pol}(D) = \operatorname{Pol}(F) \\
\varepsilon_a(F) = \varepsilon_a(F) = \text{un des deux } n\text{-gones circonscrits dans } \operatorname{Pol}(D) = \operatorname{Pol}(F)
\end{array}
\]
\note{dans la première ligne, « $F$ » est écrit sur une lettre biffée et « deux » ajouté au-dessus (lecture \uncertain{probable}) ; « circonscrits » est d'une lecture \uncertain{douteuse} ; dans la seconde, l'indice du premier $\varepsilon$ est surchargé (lu $a$) et le second membre « $\varepsilon_a(F)$ » répète le premier. La page s'arrête ici ; la suite est au lot suivant}

\end{document}
